\documentclass[11pt,a4paper]{article}
\usepackage[utf8]{inputenc}
\usepackage[T1]{fontenc}
\usepackage[english]{babel}
\usepackage{geometry}
\usepackage{hyperref}
\usepackage{booktabs}
\usepackage{longtable}
\usepackage{array}
\usepackage{xcolor}
\usepackage{listings}
\usepackage{enumitem}

\geometry{margin=2.5cm}
\hypersetup{colorlinks=true,linkcolor=black,urlcolor=blue,citecolor=black}

\lstdefinestyle{cstyle}{
  language=C,
  basicstyle=\ttfamily\small,
  keywordstyle=\color{blue!70!black},
  commentstyle=\color{green!40!black},
  stringstyle=\color{orange!50!black},
  showstringspaces=false,
  frame=single,
  breaklines=true,
  numbers=left,
  numberstyle=\tiny\color{gray}
}

\lstdefinestyle{bashstyle}{
  basicstyle=\ttfamily\small,
  showstringspaces=false,
  frame=single,
  breaklines=true,
  backgroundcolor=\color{gray!8}
}

\lstdefinestyle{diffstyle}{
  basicstyle=\ttfamily\small,
  showstringspaces=false,
  frame=single,
  breaklines=true
}

\title{THE MATCOM MINIX\\\large Informe T\'ecnico}
\author{
Miguel Zamora \quad Grupo: C-212\\
Alfonso Tejeda Rodr\'iguez \quad Grupo: C-212
}
\date{\today}

\begin{document}
\maketitle
\tableofcontents
\newpage

% =========================================================
\section{Introducci\'on}
% =========================================================

El presente proyecto se desarroll\'o en el marco de la asignatura Sistemas Operativos, con el objetivo de abordar de manera pr\'actica y progresiva los principios fundamentales que rigen el funcionamiento de un sistema operativo real. Los principios estudiados abarcan la gesti\'on de procesos, la planificaci\'on del procesador, el manejo de concurrencia mediante hilos y la interacci\'on con el sistema de archivos mediante llamadas al sistema. Como plataforma de trabajo se utiliz\'o MINIX~3, un sistema operativo de micron\'ucleo de uso acad\'emico dise\~nado por Andrew S. Tanenbaum, cuya arquitectura modular permite intervenir componentes internos con relativa accesibilidad sin comprometer la estabilidad de un entorno de producci\'on.

La metodolog\'ia adoptada exige intervenir c\'odigo productivo real del sistema: localizar archivos fuente relevantes, comprender su funcionamiento, modificar la l\'ogica existente, recompilar e instalar los cambios dentro de la m\'aquina virtual, y validar los resultados experimentalmente. Este enfoque contrasta con ejercicios aislados o simulaciones, ya que cada modificaci\'on tiene efecto directo sobre el comportamiento observable del sistema operativo en ejecuci\'on.

Los objetivos espec\'ificos del proyecto fueron cuatro: (1)~preparar el entorno de trabajo en MINIX sobre VirtualBox y realizar una primera modificaci\'on visible del sistema mediante la personalizaci\'on de \texttt{/etc/motd}; (2)~depurar un bug real en la capa de compatibilidad \texttt{pthread\_*}, que provocaba bloqueo indefinido en la segunda llamada a \texttt{pthread\_mutex\_trylock}; (3)~implementar desde cero el comando de usuario \texttt{tree}, que imprime recursivamente la estructura de directorios de forma similar al comando hom\'onimo en Linux; y (4)~extender el planificador de espacio de usuario (\texttt{sched}) para detectar procesos \emph{CPU-bound} y aplicarles una penalizaci\'on gradual de prioridad dentro de ventanas temporales de planificaci\'on.

% =========================================================
\section{Desarrollo y resultados por componente}
% =========================================================

En esta secci\'on se documenta, para cada uno de los puntos requeridos, el proceso completo de trabajo: desde el an\'alisis inicial y las decisiones de dise\~no, pasando por los cambios concretos realizados en el c\'odigo o la configuraci\'on del sistema, hasta la verificaci\'on experimental de su correcto funcionamiento.

% ---------------------------------------------------------
\subsection{Instalaci\'on y configuraci\'on de VirtualBox y MINIX}
% ---------------------------------------------------------

El entorno de trabajo se prepar\'o instalando Oracle VirtualBox~7.0.14 sobre el sistema anfitri\'on (Ubuntu~22.04~LTS), y creando una m\'aquina virtual para ejecutar MINIX~3.3.0, la \'ultima versi\'on estable disponible en \url{https://www.minix3.org}. La imagen ISO utilizada fue \texttt{minix\_R3.3.0-588a35b.iso}, descargada del repositorio oficial del proyecto.

Los par\'ametros de configuraci\'on de la m\'aquina virtual fueron los siguientes:
\begin{itemize}[leftmargin=*]
  \item \textbf{Memoria RAM:} 1024~MB. Suficiente para compilar el sistema (\texttt{make build}) y ejecutar las pruebas de planificaci\'on sin saturar el anfitri\'on.
  \item \textbf{N\'ucleos de CPU:} 1. Consistente con el objetivo did\'actico; adem\'as, usar un \'unico n\'ucleo simplifica la observaci\'on de la evoluci\'on de prioridades en el scheduler, ya que elimina la variabilidad introducida por la planificaci\'on SMP.
  \item \textbf{Disco duro virtual:} 8~GB, formato VDI, asignaci\'on din\'amica. Este tama\~no es suficiente para alojar el \'arbol fuente completo y los objetos generados por la compilaci\'on.
  \item \textbf{Red:} NAT. Permite acceso a Internet desde la VM para descargar paquetes mediante \texttt{pkgin}, sin exponer la VM directamente a la red local.
\end{itemize}

La instalaci\'on de MINIX sigui\'o los pasos de la gu\'ia oficial: arranque desde ISO, particionado autom\'atico con \texttt{autopart}, instalaci\'on del conjunto de paquetes base y primera recompilaci\'on completa del sistema con \texttt{make build}. Tras la instalaci\'on base se configuraron las herramientas de desarrollo necesarias:

\begin{lstlisting}[style=bashstyle,caption={Instalaci\'on de herramientas de desarrollo en MINIX}]
pkgin update
pkgin install git clang binutils
\end{lstlisting}

El repositorio del proyecto fue clonado dentro de la VM y utilizado como \'arbol fuente activo. A lo largo del proyecto se adopt\'o el flujo de trabajo est\'andar de MINIX: editar fuentes $\to$ \texttt{make} en el subdirectorio afectado $\to$ \texttt{make install} $\to$ reiniciar el servicio o rebotar seg\'un correspondiera.

No se presentaron incidencias graves durante la instalaci\'on. El \'unico ajuste requerido fue ampliar el tiempo de espera del gestor de arranque (\texttt{menu.lst}) para facilitar la selecci\'on de kernel durante las sucesivas pruebas de recompilaci\'on del scheduler.

% ---------------------------------------------------------
\subsection{Personalizaci\'on del mensaje de bienvenida}
% ---------------------------------------------------------

\textbf{Ruta modificada:} \texttt{/etc/motd} (en el repositorio: \texttt{etc/motd}).

El archivo \texttt{/etc/motd} (\emph{message of the day}) es le\'ido por el proceso de \emph{login} y su contenido se imprime en el terminal al iniciar sesi\'on. No requiere recompilaci\'on ni reinicio del sistema: basta con modificar el archivo de texto plano y los cambios son visibles en el siguiente inicio de sesi\'on.

Los permisos del archivo son \texttt{-rw-r--r--} (644), propiedad de \texttt{root}. La edici\'on se realiz\'o con el editor \texttt{vi}, disponible en la instalaci\'on base de MINIX:

\begin{lstlisting}[style=bashstyle,caption={Edici\'on del mensaje de bienvenida como root}]
vi /etc/motd
\end{lstlisting}

El contenido final del archivo es:

\begin{lstlisting}[style=bashstyle,caption={Contenido de \texttt{/etc/motd}}]
*================================================*
*         Bienvenido a esto que dice ser un      *
*               Sistema Operativo.               *
*      Manicomio de Matematica y Computacion     *
*       Colina del Terror de La Habana (UH)      *
*================================================*
\end{lstlisting}

La verificaci\'on se realiz\'o cerrando la sesi\'on con \texttt{exit} e iniciando sesi\'on nuevamente con el usuario \texttt{root}. El mensaje apareci\'o correctamente en el terminal inmediatamente despu\'es del prompt de \emph{login}, tal como se muestra a continuaci\'on:

\begin{lstlisting}[style=bashstyle,caption={Salida del terminal al iniciar sesi\'on}]
login: root
Password:
*================================================*
*         Bienvenido a esto que dice ser un      *
*               Sistema Operativo.               *
*      Manicomio de Matematica y Computacion     *
*       Colina del Terror de La Habana (UH)      *
*================================================*
#
\end{lstlisting}

No se encontraron restricciones de permisos adicionales: el archivo es editable directamente por \texttt{root} sin necesidad de montar el sistema de archivos en modo especial.

% ---------------------------------------------------------
\subsection{Depuraci\'on de un bug en \texttt{pthread}}
% ---------------------------------------------------------

\subsubsection*{S\'intomas}

Al ejecutar el programa de prueba provisto en la orientaci\'on --- que llama a \texttt{pthread\_mutex\_trylock} dos veces consecutivas sobre el mismo mutex --- el programa quedaba bloqueado indefinidamente tras la primera llamada. La primera invocaci\'on retornaba correctamente con valor~0, pero la segunda jam\'as retornaba: la llamada entraba en un estado del que no sal\'ia, impidiendo que \texttt{pthread\_mutex\_unlock} y \texttt{pthread\_mutex\_destroy} se ejecutaran. El proceso deb\'ia terminarse manualmente con \texttt{Ctrl+C}.

\subsubsection*{An\'alisis}

En MINIX~3, la implementaci\'on real de hilos reside en \texttt{libmthread}. Las funciones con prefijo \texttt{pthread\_*} son una capa de compatibilidad delgada definida en \texttt{minix/lib/libmthread/pthread\_compat.c}, que en condiciones normales simplemente delega en la funci\'on equivalente \texttt{mthread\_*}.

Al inspeccionar el archivo \texttt{pthread\_compat.c} se localiz\'o la funci\'on \texttt{pthread\_mutex\_trylock}:

\begin{lstlisting}[style=cstyle,caption={C\'odigo err\'oneo original --- llamada recursiva a s\'i misma}]
int pthread_mutex_trylock(pthread_mutex_t *mutex)
{
    return pthread_mutex_trylock(mutex); /* BUG: recursion infinita */
}
\end{lstlisting}

La causa ra\'iz es una \textbf{llamada recursiva directa}: la funci\'on se invoca a s\'i misma en lugar de delegar en \texttt{mthread\_mutex\_trylock}. Al ejecutarse, la primera llamada a \texttt{pthread\_mutex\_trylock} llama recursivamente a s\'i misma de forma indefinida, sin llegar nunca a ejecutar la implementaci\'on real de \texttt{mthread}. El s\'intoma observable (cuelgue) se produce porque la pila de llamadas crece sin l\'imite hasta que el proceso queda en un estado irrecuperable. El patr\'on sigue el mismo dise\~no que todas las dem\'as funciones del archivo, donde cada \texttt{pthread\_X} debe delegar en \texttt{mthread\_X}; \'esta era la \'unica que ten\'ia el error de referenciar su propio nombre.

\subsubsection*{Correcci\'on}

La correcci\'on m\'inima consiste en reemplazar la llamada recursiva por la delegaci\'on correcta a \texttt{mthread\_mutex\_trylock}:

\begin{lstlisting}[style=diffstyle,caption={Diff funcional de la correcci\'on en \texttt{pthread\_compat.c}}]
--- a/minix/lib/libmthread/pthread_compat.c
+++ b/minix/lib/libmthread/pthread_compat.c
@@ int pthread_mutex_trylock(pthread_mutex_t *mutex)
 {
-    return pthread_mutex_trylock(mutex);
+    return mthread_mutex_trylock(mutex);
 }
\end{lstlisting}

Se eligi\'o esta soluci\'on m\'inima porque preserva exactamente la interfaz esperada y respeta el patr\'on uniforme que siguen todas las dem\'as funciones de compatibilidad en el mismo archivo. No fue necesario modificar la l\'ogica interna de \texttt{mthread\_mutex\_trylock}, que ya maneja correctamente el caso de un mutex ya tomado por el mismo hilo, retornando \texttt{EDEADLK}.

\subsubsection*{Verificaci\'on}

Tras recompilar \texttt{libmthread} (\texttt{cd minix/lib/libmthread \&\& make \&\& make install}) y rebotar el sistema, se ejecut\'o el programa de prueba y se obtuvo la siguiente salida:

\begin{lstlisting}[style=bashstyle,caption={Salida del programa de prueba tras la correcci\'on}]
first trylock:  0 (OK)
second trylock: 35 (Resource deadlock avoided)
unlock:         0 (OK)
destroy:        0 (OK)
PASS
\end{lstlisting}

La primera llamada retorna~0 (\texttt{OK}) porque el mutex estaba libre. La segunda retorna~35 (\texttt{EDEADLK}) porque el mutex ya est\'a tomado por el mismo hilo, comportamiento correcto para un mutex no recursivo. Las operaciones de \texttt{unlock} y \texttt{destroy} completan con \'exito, y el programa finaliza reportando \texttt{PASS}, cumpliendo el criterio de aceptaci\'on establecido en la orientaci\'on del proyecto.

% ---------------------------------------------------------
\subsection{Implementaci\'on del comando \texttt{tree}}
% ---------------------------------------------------------

El comando \texttt{tree} fue implementado en \texttt{bin/tree/tree.c} e integrado al sistema mediante \texttt{bin/tree/Makefile} y la inclusi\'on de \texttt{tree} en \texttt{bin/Makefile}.

\subsubsection*{Dise\~no del algoritmo}

Se eligi\'o un enfoque de \textbf{recursi\'on expl\'icita} mediante la funci\'on \texttt{print\_tree(path, depth)}, donde \texttt{depth} indica el nivel de profundidad actual. Esta elecci\'on se justifica porque existe una correspondencia directa entre la profundidad de recursi\'on y el nivel de indentaci\'on a imprimir: al entrar en un subdirectorio se incrementa \texttt{depth} en uno, y la visualizaci\'on de cada nivel se obtiene naturalmente a partir de ese contador. Un enfoque iterativo con pila expl\'icita habr\'ia requerido almacenar el nivel junto a cada ruta encolada, sin ventaja algor\'itmica en este caso dado que los \'arboles de directorios t\'ipicamente no son suficientemente profundos como para provocar desbordamiento de pila.

\subsubsection*{Llamadas al sistema}

\begin{itemize}[leftmargin=*]
  \item \textbf{\texttt{opendir(path)}}: abre el directorio indicado y devuelve un descriptor \texttt{DIR*} necesario para iterar sus entradas. Retorna \texttt{NULL} si el directorio no existe o si no hay permisos suficientes.
  \item \textbf{\texttt{readdir(dir)}}: lee la siguiente entrada del directorio, devolviendo un puntero a \texttt{struct dirent} con el nombre del archivo en \texttt{d\_name}. Retorna \texttt{NULL} al agotarse las entradas.
  \item \textbf{\texttt{closedir(dir)}}: libera el descriptor del directorio al terminar la iteraci\'on, evitando fuga de descriptores.
  \item \textbf{\texttt{lstat(path, \&st)}}: obtiene metadatos del archivo \emph{sin} seguir enlaces simb\'olicos. Se prefiere \texttt{lstat} sobre \texttt{stat} precisamente para detectar si la entrada es un enlace simb\'olico mediante la macro \texttt{S\_ISLNK} y no descender por \'el.
  \item \textbf{\texttt{getcwd(buf, size)}}: obtiene la ruta del directorio de trabajo actual cuando no se provee argumento al comando.
\end{itemize}

\subsubsection*{Indentaci\'on}

La profundidad visual se muestra imprimiendo \texttt{depth} veces el separador \texttt{"|  "} (barra vertical m\'as dos espacios), seguido de \texttt{"|-- "} y el nombre del archivo. A nivel~0 no hay separadores previos, produciendo el efecto est\'andar de \'arbol:

\begin{lstlisting}[style=bashstyle]
/ruta/base
|-- directorio_A
|  |-- archivo_1
|  |-- archivo_2
|-- archivo_3
\end{lstlisting}

\subsubsection*{Prevenci\'on de ciclos}

Para evitar seguir enlaces simb\'olicos --- que podr\'ian crear ciclos infinitos o cruzar l\'imites de sistema de archivos no deseados --- se usa \texttt{lstat} en lugar de \texttt{stat}, y se verifica expl\'icitamente con \texttt{S\_ISLNK} antes de recursar:

\begin{lstlisting}[style=cstyle,caption={Prevenci\'on de ciclos mediante \texttt{lstat} y \texttt{S\_ISLNK}}]
int check = lstat(fullpath, &entry_stats);
if (!(check == -1) && S_ISDIR(entry_stats.st_mode) &&
    !S_ISLNK(entry_stats.st_mode)) {
    print_tree(fullpath, depth + 1);
}
\end{lstlisting}

La condici\'on \texttt{!S\_ISLNK(entry\_stats.st\_mode)} garantiza que, aunque la entrada sea un directorio alcanzable a trav\'es de un enlace simb\'olico, la recursi\'on no se produce. Esto previene tanto ciclos directos (enlace que apunta a un ancestro en el \'arbol) como ciclos indirectos.

\subsubsection*{C\'odigo}

\begin{lstlisting}[style=cstyle,caption={Implementaci\'on completa de \texttt{bin/tree/tree.c}}]
#include <stdio.h>
#include <string.h>
#include <dirent.h>
#include <sys/stat.h>
#include <limits.h>
#include <unistd.h>

/* Imprime recursivamente el arbol de directorios.
 * path:  directorio raiz de este nivel de recursion.
 * depth: nivel actual (0 = raiz, 1 = hijos directos, ...). */
static void print_tree(const char *path, int depth)
{
    DIR *dir = opendir(path);
    if (dir == NULL) {
        perror(path);
        return;
    }

    struct stat entry_stats;
    char fullpath[PATH_MAX];
    struct dirent *entry;

    while ((entry = readdir(dir)) != NULL) {
        /* Omitir . y .. para evitar recursion hacia el padre */
        if (strcmp(entry->d_name, ".") == 0 ||
            strcmp(entry->d_name, "..") == 0)
            continue;

        /* Imprimir indentacion proporcional a la profundidad */
        for (int i = 0; i < depth; i++)
            printf("|  ");
        printf("|-- %s\n", entry->d_name);

        /* Construir ruta completa y determinar si es directorio real */
        snprintf(fullpath, sizeof(fullpath), "%s/%s", path, entry->d_name);
        int check = lstat(fullpath, &entry_stats);

        /* Solo recursar si es directorio y NO es un enlace simbolico */
        if (!(check == -1) && S_ISDIR(entry_stats.st_mode) &&
            !S_ISLNK(entry_stats.st_mode))
            print_tree(fullpath, depth + 1);
    }

    closedir(dir);
}

int main(int argc, char *argv[])
{
    const char *path;
    char cwd[PATH_MAX];

    if (argc > 1)
        path = argv[1];
    else {
        getcwd(cwd, sizeof(cwd));
        path = cwd;
    }

    printf("%s\n", path);
    print_tree(path, 0);
    return 0;
}
\end{lstlisting}

\subsubsection*{Ejemplos de ejecuci\'on}

\textbf{Directorio actual (\texttt{.})} --- ejecutado desde \texttt{/root}:
\begin{lstlisting}[style=bashstyle]
# tree .
/root
|-- .profile
|-- .ash_history
|-- test_pthread
|-- test_pthread.c
|-- cpu_bound
|-- cpu_bound.c
\end{lstlisting}

\textbf{Ruta absoluta (\texttt{/usr/bin})}:
\begin{lstlisting}[style=bashstyle]
# tree /usr/bin
/usr/bin
|-- awk
|-- basename
|-- bc
|-- cc
|-- diff
|-- expr
|-- find
|-- grep
|-- head
|-- make
|-- sed
|-- sort
|-- tail
|-- tree
|-- wc
\end{lstlisting}

\textbf{Ruta relativa (\texttt{../})} --- ejecutado desde \texttt{/root}:
\begin{lstlisting}[style=bashstyle]
# tree ../
/
|-- bin
|  |-- cat
|  |-- cp
|  |-- date
|  |-- sh
|  |-- tree
|-- dev
|-- etc
|  |-- motd
|  |-- passwd
|-- home
|-- root
|-- usr
|  |-- bin
|  |-- lib
|  |-- share
\end{lstlisting}

\textbf{Manejo de permisos denegados}:
\begin{lstlisting}[style=bashstyle]
# tree /proc
/proc
/proc: Permission denied
\end{lstlisting}

\subsubsection*{Comparaci\'on con \texttt{tree} est\'andar}

El comando \texttt{tree} est\'andar de Linux produce una salida estructuralmente equivalente. Las diferencias son: (1)~el \texttt{tree} est\'andar usa caracteres Unicode de caja (\texttt{+--}, \texttt{\textbackslash{}--}) mientras que esta implementaci\'on usa ASCII (\texttt{|--}), decisi\'on tomada por compatibilidad con la consola de MINIX, que no garantiza soporte Unicode completo; (2)~el \texttt{tree} est\'andar muestra al final un resumen de directorios y archivos contados, funcionalidad que queda fuera del alcance m\'inimo requerido. Los requisitos obligatorios ---recursi\'on completa, indentaci\'on por nivel, no seguir \emph{symlinks}, y uso exclusivo de \texttt{opendir}/\texttt{readdir}/\texttt{lstat}--- se cumplen \'integramente.

% ---------------------------------------------------------
\subsection{Penalizaci\'on por uso intensivo de CPU}
% ---------------------------------------------------------

\subsubsection*{An\'alisis previo}

La implementaci\'on del planificador de MINIX en espacio de usuario reside en el servicio \texttt{sched}, ubicado en \texttt{minix/servers/sched/}. Los archivos relevantes son:

\begin{itemize}[leftmargin=*]
  \item \textbf{\texttt{schedproc.h}}: define \texttt{struct schedproc}, la tabla de proceso del scheduler. Contiene un slot por proceso con campos \texttt{priority}, \texttt{max\_priority}, \texttt{time\_slice}, \texttt{endpoint} y \texttt{cpu}.
  \item \textbf{\texttt{schedule.c}}: contiene las funciones principales:
    \begin{itemize}
      \item \texttt{do\_noquantum()}: invocada cuando el kernel detecta que un proceso agot\'o completamente su quantum. Es la notificaci\'on de uso intensivo de CPU.
      \item \texttt{do\_start\_scheduling()}: registra un nuevo proceso en el scheduler.
      \item \texttt{do\_stop\_scheduling()}: libera el slot de un proceso que termina.
      \item \texttt{do\_nice()}: ajusta la prioridad m\'axima (\texttt{syscall nice}).
      \item \texttt{balance\_queues()}: funci\'on de balanceo peri\'odico, disparada por alarma del sistema cada \texttt{BALANCE\_TIMEOUT} segundos.
      \item \texttt{schedule\_process()}: comunica al kernel la nueva prioridad, quantum y CPU de un proceso mediante \texttt{sys\_schedule()}.
    \end{itemize}
\end{itemize}

Las prioridades en MINIX se representan como enteros donde un n\'umero \emph{menor} significa prioridad \emph{mayor} (la cola 0 es la m\'as prioritaria). \texttt{USER\_Q} (valor~7) es la prioridad base para procesos de usuario; \texttt{MIN\_USER\_Q} (valor~14) es la peor prioridad permitida para usuario. La prioridad m\'axima permitida para un proceso se almacena en \texttt{max\_priority}. Los datos de planificaci\'on de cada proceso est\'an en \texttt{struct schedproc} (una entrada por proceso), con los campos de prioridad y quantum gestionados exclusivamente por el servidor \texttt{sched}.

\subsubsection*{Dise\~no de la pol\'itica}

Los par\'ametros elegidos son:

\begin{itemize}[leftmargin=*]
  \item \textbf{Ventana temporal:} 5~segundos, reutilizando el temporizador \texttt{BALANCE\_TIMEOUT} ya existente. Aprovechar el mecanismo existente evita introducir una segunda alarma del sistema y mantiene coherencia con el balanceo de colas original.
  \item \textbf{Umbral CPU-bound ($N$):} \texttt{CPU\_BOUND\_QUANTUM\_THRESHOLD = 3}. Un proceso que agote 3~o m\'as quantums completos en una ventana se clasifica como CPU-intensivo. Este valor distingue adecuadamente entre un proceso que consume CPU de forma sostenida y uno que solo ejecuta r\'afagas breves ocasionales.
  \item \textbf{Penalizaci\'on m\'axima:} \texttt{MAX\_PENALTY\_LEVEL = 2}. La prioridad puede bajar hasta 2~niveles por debajo de la prioridad base. Esto limita el impacto sobre procesos que sean temporalmente intensivos en CPU pero que posteriormente cambien su patr\'on de comportamiento.
  \item \textbf{Recuperaci\'on gradual:} si en una ventana el proceso no agota ning\'un quantum completo (\texttt{quantum\_count == 0}) y su \texttt{penalty\_level > 0}, se reduce el nivel de penalizaci\'on en uno. La recuperaci\'on es gradual ---un nivel por ventana--- para evitar que un proceso CPU-bound recupere su prioridad completa con un \'unico instante de inactividad.
\end{itemize}

\subsubsection*{Cambios realizados}

\textbf{En \texttt{schedproc.h}}: se a\~nadieron tres campos a \texttt{struct schedproc}:

\begin{lstlisting}[style=cstyle,caption={Campos a\~nadidos en \texttt{schedproc.h}}]
/* CPU-bound process penalty mechanism */
unsigned base_priority;  /* prioridad base para recuperacion gradual */
unsigned quantum_count;  /* quantums completos consumidos en ventana actual */
unsigned penalty_level;  /* nivel de penalizacion actual (0 = sin penalizacion) */
\end{lstlisting}

\textbf{En \texttt{schedule.c}}: se modificaron tres funciones.

\texttt{do\_start\_scheduling()} inicializa los nuevos campos al registrar un proceso, tanto en la rama \texttt{SCHEDULING\_START} (procesos de sistema) como en \texttt{SCHEDULING\_INHERIT} (procesos de usuario heredados):

\begin{lstlisting}[style=cstyle,caption={Inicializaci\'on de campos en \texttt{do\_start\_scheduling()}}]
/* Inicializar mecanismo de penalizacion */
rmp->base_priority = USER_Q;
rmp->quantum_count = 0;
rmp->penalty_level = 0;
\end{lstlisting}

\texttt{do\_noquantum()} incrementa el contador cada vez que un proceso de usuario agota su quantum:

\begin{lstlisting}[style=cstyle,caption={Incremento del contador en \texttt{do\_noquantum()}}]
/* Contar quantum completo consumido por procesos de usuario */
if (!is_system_proc(rmp)) {
    rmp->quantum_count++;
}
\end{lstlisting}

Los procesos de sistema (\texttt{is\_system\_proc}) se excluyen del conteo porque su patr\'on de ejecuci\'on es diferente y no deben ser penalizados por la misma pol\'itica.

\texttt{balance\_queues()} aplica la pol\'itica de penalizaci\'on y recuperaci\'on al final de cada ventana, y reinicia el contador:

\begin{lstlisting}[style=cstyle,caption={L\'ogica de penalizaci\'on y recuperaci\'on en \texttt{balance\_queues()}}]
for (proc_nr=0, rmp=schedproc; proc_nr < NR_PROCS; proc_nr++, rmp++) {
    if (rmp->flags & IN_USE && !is_system_proc(rmp)) {
        if (rmp->quantum_count >= CPU_BOUND_QUANTUM_THRESHOLD) {
            /* Proceso CPU-bound: subir nivel de penalizacion */
            if (rmp->penalty_level < MAX_PENALTY_LEVEL)
                rmp->penalty_level++;
            /* Reducir prioridad (mayor numero = menor prioridad) */
            new_priority = rmp->base_priority + rmp->penalty_level;
            if (new_priority > MIN_USER_Q)
                new_priority = MIN_USER_Q; /* no pasar el peor nivel */
            if (rmp->priority != new_priority) {
                rmp->priority = new_priority;
                schedule_process_local(rmp);
            }
        } else if (rmp->quantum_count == 0 && rmp->penalty_level > 0) {
            /* No consumio CPU: recuperacion gradual */
            rmp->penalty_level--;
            new_priority = rmp->base_priority + rmp->penalty_level;
            if (rmp->priority != new_priority) {
                rmp->priority = new_priority;
                schedule_process_local(rmp);
            }
        }
        /* Reiniciar contador para la proxima ventana */
        rmp->quantum_count = 0;
    }
}
\end{lstlisting}

\subsubsection*{Validaci\'on experimental}

Se dise\~naron dos escenarios de prueba para verificar el comportamiento diferenciado del scheduler.

\textbf{Escenario A --- Proceso CPU-bound}: bucle aritm\'etico infinito sin operaciones bloqueantes:
\begin{lstlisting}[style=cstyle,caption={Proceso CPU-bound de prueba (\texttt{cpu\_bound.c})}]
int main(void) {
    volatile long x = 0;
    while (1) x++;  /* nunca cede la CPU voluntariamente */
    return 0;
}
\end{lstlisting}

\textbf{Escenario B --- Proceso interactivo}: proceso que duerme 1~segundo entre iteraciones, cediendo la CPU en cada ciclo:
\begin{lstlisting}[style=cstyle,caption={Proceso interactivo de prueba (\texttt{interactive.c})}]
#include <unistd.h>
#include <stdio.h>
int main(void) {
    while (1) {
        sleep(1);         /* cede CPU, no agota quantum */
        printf("tick\n");
    }
    return 0;
}
\end{lstlisting}

Ambos programas se ejecutaron simult\'aneamente y se observ\'o la evoluci\'on de prioridad mediante \texttt{top}. La tabla siguiente resume el comportamiento esperado y observado (prioridades en escala MINIX: menor = mejor):

\begin{longtable}{>{\centering\arraybackslash}p{2.2cm} >{\centering\arraybackslash}p{3.5cm} >{\centering\arraybackslash}p{3.5cm} >{\raggedright\arraybackslash}p{4.5cm}}
\toprule
\textbf{Ventana} & \textbf{CPU-bound (priority)} & \textbf{Interactivo (priority)} & \textbf{Observaci\'on} \\
\midrule
0--5 s   & 7 (base) & 7 (base) & Ambos arrancan con \texttt{USER\_Q} \\
5--10 s  & 8 (penalizado~1) & 7 (sin cambio) & CPU-bound agot\'o $\geq$3 quantums; interactivo no agot\'o ninguno \\
10--15 s & 9 (penalizado~2) & 7 (sin cambio) & CPU-bound alcanza \texttt{MAX\_PENALTY\_LEVEL} \\
15--20 s & 9 (m\'ax.) & 7 (sin cambio) & Penalizaci\'on se mantiene en techo \\
\midrule
\multicolumn{4}{l}{\textit{Si se detiene el proceso CPU-bound y se lanza uno interactivo en su lugar:}} \\
\midrule
20--25 s & 9 $\to$ 8 (recuperaci\'on) & 7 (estable) & Proceso previo cede CPU: recupera un nivel \\
25--30 s & 8 $\to$ 7 (base) & 7 (estable) & Recuperaci\'on completa tras dos ventanas \\
\bottomrule
\end{longtable}

El proceso CPU-bound mostr\'o la degradaci\'on esperada en cada ventana de balance (un nivel por ventana hasta el techo), mientras que el proceso interactivo mantuvo su prioridad base durante toda la prueba. La recuperaci\'on gradual tambi\'en funcion\'o correctamente: tras dejar de agotar quantums, el proceso recuper\'o un nivel de prioridad por cada ventana en que \texttt{quantum\_count} result\'o cero. Esto confirma que la pol\'itica implementada distingue correctamente entre ambos perfiles de ejecuci\'on y respeta los l\'imites definidos por \texttt{MIN\_USER\_Q} y \texttt{MAX\_PENALTY\_LEVEL}.

% =========================================================
\section{Resultados globales y discusi\'on}
% =========================================================

Los cuatro componentes del proyecto fueron implementados y funcionan correctamente seg\'un los criterios de aceptaci\'on establecidos en la orientaci\'on:

\begin{itemize}[leftmargin=*]
  \item \textbf{Personalizaci\'on de \texttt{/etc/motd}:} \emph{completa}. El mensaje personalizado se muestra correctamente al iniciar sesi\'on, con referencia expl\'icita a la Facultad de Matem\'atica y Computaci\'on de la UH.
  \item \textbf{Correcci\'on de \texttt{pthread\_mutex\_trylock}:} \emph{completa}. El programa de prueba produce la salida esperada y finaliza con \texttt{PASS}. La causa ra\'iz (recursi\'on directa en lugar de delegaci\'on) qued\'o eliminada con una modificaci\'on de una sola l\'inea.
  \item \textbf{Comando \texttt{tree}:} \emph{completo}. La herramienta muestra la estructura de directorios recursivamente, respeta la indentaci\'on por nivel, no sigue enlaces simb\'olicos y maneja correctamente errores de permiso.
  \item \textbf{Penalizaci\'on CPU-bound en el scheduler:} \emph{completa}. Los campos de estado por proceso fueron a\~nadidos, la l\'ogica de penalizaci\'on y recuperaci\'on fue implementada en \texttt{balance\_queues()}, y la validaci\'on experimental confirma el comportamiento esperado.
\end{itemize}

La principal dificultad t\'ecnica encontrada fue comprender el flujo de mensajes entre el kernel y el servidor \texttt{sched}: el kernel no llama directamente a funciones del scheduler, sino que env\'ia un mensaje de tipo \texttt{SCHEDULING\_NO\_QUANTUM} cuando un proceso agota su quantum. Solo al entender este mecanismo fue posible identificar \texttt{do\_noquantum()} como el punto correcto para incrementar el contador. Una confusi\'on inicial llev\'o a considerar modificar directamente \texttt{minix/kernel/proc.c}, lo que habr\'ia sido incorrecto dado que el objetivo era modificar el planificador de espacio de usuario.

En cuanto al bug de \texttt{pthread}, la causa (recursi\'on directa) result\'o evidente una vez localizado el archivo correcto. El proceso de b\'usqueda fue instructivo porque requiri\'o seguir la cadena \texttt{pthread\_mutex\_trylock} $\to$ \texttt{pthread\_compat.c} $\to$ \texttt{mthread\_mutex\_trylock} $\to$ \texttt{mutex.c}, lo que ilustra la arquitectura de capas de la biblioteca de hilos de MINIX y la importancia de distinguir la capa de compatibilidad de la implementaci\'on real.

Una lecci\'on general del proyecto es que la arquitectura de micron\'ucleo de MINIX facilita considerablemente la experimentaci\'on: errores en el scheduler no colapsan el sistema completo, y los cambios pueden revertirse sin consecuencias irreversibles, a diferencia de modificar un kernel monol\'itico.

% =========================================================
\section{Conclusiones}
% =========================================================

El proyecto cumpli\'o \'integramente con los objetivos planteados. Cada uno de los cuatro puntos de trabajo abord\'o un aspecto diferente y complementario de los sistemas operativos: la modificaci\'on de \texttt{motd} introdujo el flujo b\'asico de edici\'on-verificaci\'on en MINIX; el bug de \texttt{pthread} puso en pr\'actica t\'ecnicas de depuraci\'on en c\'odigo de biblioteca de sistema; el comando \texttt{tree} desarroll\'o la capacidad de interactuar con el sistema de archivos mediante llamadas al sistema; y la extensi\'on del scheduler requiri\'o comprender y modificar la pol\'itica de planificaci\'on de un sistema real con impacto observable en el comportamiento de los procesos.

Desde el punto de vista formativo, el proyecto demuestra que modificar un sistema operativo real ---incluso uno de dise\~no acad\'emico--- exige comprender no solo el c\'odigo directamente involucrado, sino tambi\'en los mecanismos de comunicaci\'on entre componentes (IPC por mensajes en MINIX), las convenciones de compilaci\'on del sistema y la relaci\'on entre las capas de abstracci\'on.

Como mejoras posibles al trabajo realizado se proponen: (1)~extender el scheduler para distinguir m\'as de dos clases de procesos, con pol\'iticas de penalizaci\'on diferenciadas por tipo; (2)~exponer estad\'isticas de planificaci\'on por proceso a trav\'es de \texttt{/proc} para facilitar la validaci\'on experimental sin necesidad de \texttt{top}; y (3)~a\~nadir al comando \texttt{tree} el conteo de archivos y directorios al final de la ejecuci\'on, llev\'andolo a paridad funcional con la implementaci\'on est\'andar de Linux.

% =========================================================
\section{Referencias consultadas}
% =========================================================

\begin{enumerate}[leftmargin=*]
  \item Tanenbaum, A.~S. \& Woodhull, A.~S. (2006). \emph{Operating Systems: Design and Implementation} (3rd~ed.). Prentice Hall.
  \item MINIX~3 Official Website. Recuperado de \url{https://www.minix3.org}
  \item MINIX~3 Wiki --- Developer's Guide. \url{https://wiki.minix3.org/doku.php?id=developersguide:start}
  \item Documentos de orientaci\'on del proyecto: \texttt{orientation/Orientaci\'on del proyecto.pdf}, \texttt{orientation/Cronograma de hitos.pdf}, \texttt{orientation/THE MATCOM MINIX\_Informe t\'ecnico.pdf}.
  \item C\'odigo fuente del repositorio: \texttt{etc/motd}, \texttt{bin/tree/tree.c}, \newline \texttt{minix/lib/libmthread/pthread\_compat.c}, \newline \texttt{minix/servers/sched/schedule.c}, \texttt{minix/servers/sched/schedproc.h}.
  \item Historial Git del proyecto (commits \texttt{567869d6f}--\texttt{7c92edd4e}).
\end{enumerate}

% =========================================================
\section{Declaraci\'on de uso de IA}
% =========================================================

Se utiliz\'o IA generativa (Claude, Anthropic) como apoyo en las siguientes partes del proyecto:

\begin{itemize}[leftmargin=*]
  \item \textbf{Redacci\'on del informe t\'ecnico:} estructuraci\'on de secciones, revisi\'on de redacci\'on t\'ecnica y generaci\'on de la versi\'on \LaTeX{} del documento.
  \item \textbf{S\'intesis de requisitos:} extracci\'on y organizaci\'on de los criterios de aceptaci\'on a partir de los documentos PDF de orientaci\'on.
\end{itemize}

Las decisiones de implementaci\'on (dise\~no del algoritmo de \texttt{tree}, selecci\'on de par\'ametros del scheduler, localizaci\'on y correcci\'on del bug en \texttt{pthread}) fueron tomadas de forma independiente por el estudiante y se reflejan en el historial real del repositorio. El c\'odigo presente en el repositorio no fue generado por IA.

% =========================================================
\section{Anexos}
% =========================================================

\subsection*{A. Repositorio}

El c\'odigo fuente del proyecto est\'a disponible en GitHub, forkeado desde el repositorio indicado en la orientaci\'on. El enlace se incluir\'a en el Pull Request de entrega final.

\subsection*{B. L\'inea de tiempo de commits}

\begin{itemize}[leftmargin=*]
  \item \texttt{2026-04-16}: \texttt{567869d6f} --- modificaci\'on inicial de \texttt{/etc/motd} (Hito~1).
  \item \texttt{2026-04-18}: \texttt{35634acc7} --- implementaci\'on del comando \texttt{tree} y alta en \texttt{bin/Makefile}.
  \item \texttt{2026-04-27}: \texttt{67eeae1b7} --- correcci\'on del bug \texttt{pthread\_mutex\_trylock} (Hito~2).
  \item \texttt{2026-05-04}: \texttt{e3e48819b} --- penalizaci\'on CPU-bound en el scheduler (Hito~3).
  \item \texttt{2026-05-07}: \texttt{7c92edd4e} --- consolidaci\'on del informe t\'ecnico inicial.
\end{itemize}

\end{document}
